\chapter{The Bond-Futures Basis Trade}\label{ch:s2:the-bond-futures-basis-trade}

Buy a Treasury, sell the future on it, finance the bond in repo, and earn the few basis points by which the future's implied repo rate exceeds the rate
you pay. At its peak before March 2020, the Office of Financial Research estimated hedge funds' basis-trade positions at \$400 to \$500 billion. That
was more than 60\% of their Treasury exposure and more than 70\% of their repo borrowing. Between 18 February and 17 March 2020 they cut their short
positions in the 2-, 5- and 10-year futures from \$659 billion to \$554 billion. The Federal Reserve then committed to buy Treasuries in whatever
amounts were needed. On this chapter's synthetic market a basis book at 50 times its capital earns 10.3\% a year until a planted ten-day funding
stress. Margin calls force it to sell at the worst basis, and it ends the stress's aftermath 9.6\% poorer than it began. At ten times, it earns 2.1\%
and ends only 0.2\% down. The build is \texttt{firm.basistrade}.

\section{Mechanics: cash, future, repo}

A Treasury future (Book~2, chapter~6) is delivered by the short, who chooses the cheapest-to-deliver bond from a basket. The gross basis is the bond's
price minus the futures price times its conversion factor. Net of the carry to delivery, it is the net basis. The implied repo rate is the financing
rate at which buying the bond, delivering it into the future and repaying the loan breaks even. When the implied repo rate exceeds the repo rate, a
long basis position earns the difference: long the bond financed in repo, short the future. The book below earns 20 basis points a year on the face it
holds, the chapter's own assumption.

The trade has three legs of risk. The basis can move against the book before delivery. The repo loan must be rolled, often overnight. And the futures
exchange and the repo lender each hold collateral: a margin on the futures and a haircut on the bond.

\begin{definition}[Basis-trade leverage]\label{def:s2:the-bond-futures-basis-trade:leverage}
\emph{Basis-trade leverage}\index{basis-trade leverage} is the face value of bonds held long (and futures short) per unit of the trader's capital; it
is capped by the collateral the trade needs, so at most one divided by the repo haircut plus the futures margin (per unit of face), and it falls when
either rises.
\end{definition}

\begin{definition}[Repo rollover risk]\label{def:s2:the-bond-futures-basis-trade:rollover}
\emph{Repo rollover risk}\index{repo rollover risk} is the risk that a position financed in short-term repo cannot be refinanced on the same terms
when the loan matures: at a higher rate, a higher haircut, a smaller amount, or not at all.
\end{definition}

\section{Leverage and who does it}

Barth and Kahn traced the trade in regulatory data. Asset managers wanted duration without using their balance sheets, and bought Treasury futures.
That pushed futures above their no-arbitrage price, and hedge funds met the demand by going long the basis. From the end of 2017 to September 2019,
hedge funds' Treasury exposure grew from \$1.06 trillion to \$2.02 trillion and their short futures by \$352 billion. The authors describe the funds as
warehouses. They store Treasuries for the holders of long futures and fund them in repo, and in doing so link the Treasury, futures and repo markets.

\begin{dated}{2026-09}{Leveraged funds in Treasury futures}\label{dat:s2:the-bond-futures-basis-trade:tff}
The CFTC's Traders in Financial Futures report places each large trader in a category; ``leveraged funds'' include hedge funds and other managed
money, and hold positions of many kinds, not only basis trades. Summing their net positions in the six CBOT Treasury futures (2-, 5-, 10-year, Ultra
10-year, bond and Ultra bond) at face value, they were net long \$34 billion on average from July 2010 to 2014 and net short \$116 billion at the end of
2017. They were net short \$475 billion on 18 February 2020 and \$376 billion on 17 March 2020, reached a net short of \$1,184 billion on 12 November
2024, and were net short \$787 billion on 15 September 2026.
\end{dated}

The public series (\cref{fig:s2:the-bond-futures-basis-trade:tff}) shows the crowding the trade depends on, and the size of the flow that would come if
it unwound.

\begin{omfigure}
\begin{tikzpicture}
\begin{axis}[width=11cm,height=5.4cm,xlabel={\small year},ylabel={\small net position (\$ billion)},tick label style={font=\footnotesize},
  xmin=2010.5,xmax=2027,ymin=-1250,ymax=250,grid=major,grid style={gray!25},table/col sep=comma,/pgf/number format/1000 sep={}]
\addplot[omThm,thick] table[x=year,y=total]{figdata/strategies-2/11-the-bond-futures-basis-trade/tff.csv};
\addplot[black,thin,domain=2010.5:2027] {0};
\end{axis}
\end{tikzpicture}

\omcaption{Leveraged funds' weekly net position in the six CBOT Treasury futures, at face value, July 2010 to September 2026 (negative: net short).
The category holds more than basis trades. Derived from the CFTC's Traders in Financial Futures files. Data:
\texttt{s2\_fetch\_tff}.}\label{fig:s2:the-bond-futures-basis-trade:tff}
\end{omfigure}

\section{March 2020}

Barth and Kahn describe the stress. In the first week of March 2020, bid--ask spreads on longer Treasuries widened, repo rates on Treasury collateral
jumped, and arbitrage spreads diverged. Sales by real-money investors raised futures margins and made repo volatile. Hedge funds partly unwound their
basis positions: their short 2-, 5- and 10-year futures fell from \$659 billion to \$554 billion between 18 February and 17 March. The authors estimate
that large basis traders sold \$91 to \$105 billion of Treasuries between the end of February and the end of March. They also argue that the basis
trade was unlikely to be the primary cause of the stress before 17 March, since the Treasuries it held kept trading at a premium. On 15 March the
Federal Open Market Committee said it would increase its holdings of Treasuries by at least \$500 billion. On 23 March it said it would buy them in the
amounts needed to support smooth market functioning.

The synthetic stress has the same ingredients. For ten days from year 2.5 the basis moves 60 cents per 100 of face against the book. The repo haircut
and the futures margin rise from 1\% to 3\% each, and the repo rate rises by one percentage point. Afterwards the basis comes back over two months,
while the haircut and margin stay high. A book whose capital no longer covers its collateral must sell down to what its capital funds, paying 10 cents
per 100 on what it sells, and cannot re-lever until the stress has passed (\cref{lst:s2:the-bond-futures-basis-trade:book}).

\begin{center}
\small
\begin{tabular}{@{}lcccc@{}}
\toprule
five years, capital 1 & 10$\times$ & 20$\times$ & 30$\times$ & 50$\times$\\
\midrule
annual return before the stress (\%) & 2.1 & 4.1 & 6.2 & 10.3\\
trough in the stress (\%) & $-6.1$ & $-10.8$ & $-13.2$ & $-17.8$\\
after the basis recovered (\%) & $-0.2$ & $-1.9$ & $-4.5$ & $-9.6$\\
days with forced sales & 0 & 10 & 10 & 66\\
face held after the stress & 10.4 & 16.5 & 16.8 & 17.5\\
five-year annual return (\%) & 1.9 & 3.5 & 4.8 & 7.1\\
\bottomrule
\end{tabular}
\end{center}

At ten times, the book never falls short of its collateral. It rides out the move and recovers almost all of it. At fifty times, the book already sits at
its collateral cap, one over 2\%, so even ordinary days force small sales: 25 before the stress. In the stress the higher haircut and margin cap its
leverage at 16.7, and its face falls from 60.7 to 17.5, sold into the stressed basis. Over the stress and recovery its P\&L as a share of capital was $-6.6\%$ from marks, $-3.3\%$ from forced-sale costs, $-0.6\%$ from
dearer repo, and $+0.8\%$ from the spread still earned. The marks came back for the positions it kept, not for the ones it sold
(\cref{fig:s2:the-bond-futures-basis-trade:capital}).

\begin{omfigure}
\begin{tikzpicture}
\begin{axis}[width=11cm,height=5.4cm,xlabel={\small year},ylabel={\small capital (start 1)},tick label style={font=\footnotesize},
  xmin=0,xmax=5,ymin=0.95,ymax=1.45,grid=major,grid style={gray!25},table/col sep=comma,
  legend style={font=\scriptsize,at={(0.03,0.97)},anchor=north west}]
\addplot[omDef,thick] table[x=year,y=L10]{figdata/strategies-2/11-the-bond-futures-basis-trade/capital.csv};
\addplot[gray,thick] table[x=year,y=L20]{figdata/strategies-2/11-the-bond-futures-basis-trade/capital.csv};
\addplot[omThm,thick] table[x=year,y=L50]{figdata/strategies-2/11-the-bond-futures-basis-trade/capital.csv};
\legend{10$\times$,20$\times$,50$\times$}
\end{axis}
\end{tikzpicture}

\omcaption{Capital of the synthetic basis book at three leverages over five years, with a funding stress at year 2.5. The fifty-times book's
losses are locked in by forced sales; the ten-times book rides the basis back. Data: \texttt{s2\_basis.paths}.}\label{fig:s2:the-bond-futures-basis-trade:capital}
\end{omfigure}

\section{Measuring the crowding}

The trade's risk is not in any one book: it is the sum of everyone's. Public measures include leveraged funds' short futures in the CFTC report, dealers'
repo lending, and the gap between futures-implied and cash yields. None of them separates basis trades from other positions. Barth and Kahn used
regulatory filings that are not public. A trader running the book should measure what the market would have to absorb if every holder's collateral
rose at once. The synthetic stress shows how quickly that collateral cap binds when it does.

\section{Strategy files}

\begin{strategyfile}{Cash-futures basis long}\label{strat:s2:the-bond-futures-basis-trade:long}
\sfield{Who pays you, and why} Asset managers who want duration through futures without using their balance sheets, and push futures above
fair value.
\sfield{Instruments and venues} Treasury notes bought and financed in repo; Treasury futures sold on CBOT.
\sfield{Signal} Implied repo rate against the repo rate for the cheapest-to-deliver.
\sfield{Sizing and execution} Leverage well inside one over (haircut plus margin) under stressed terms; repo term matched to the trade.
\sfield{Costs} Repo spread, futures fees, roll costs.
\sfield{How it dies} A rise in haircut and margin that forces sales into a moving basis; repo that cannot be rolled.
\sfield{Horizon, capacity, infrastructure} Quarterly rolls; repo and futures clearing lines.
\sfield{Backtest honestly} Stress collateral terms, not only prices.
\sfield{Sources} Barth and Kahn (2021): \$400--500 billion at the peak; this chapter: 10.3\% a year at 50 times and $-9.6\%$ after a funding stress.
\end{strategyfile}

\begin{strategyfile}{Delivery option trade}\label{strat:s2:the-bond-futures-basis-trade:delivery}
\sfield{Who pays you, and why} Futures longs who give the short the choice of bond and timing.
\sfield{Instruments and venues} Deliverable bonds; the future.
\sfield{Signal} The net basis against the value of the switch option as rates and the curve move.
\sfield{Sizing and execution} Long the net basis of bonds likely to become cheapest.
\sfield{Costs} Bid-ask on off-the-run bonds.
\sfield{How it dies} Rates stay where the cheapest bond does not change.
\sfield{Horizon, capacity, infrastructure} A delivery cycle; a delivery-option model.
\sfield{Backtest honestly} Actual delivery baskets and conversion factors.
\sfield{Sources} No performance figure verified.
\end{strategyfile}

\begin{strategyfile}{Basis-trade crowding signal}\label{strat:s2:the-bond-futures-basis-trade:crowding}
\sfield{Who pays you, and why} Forced sellers in an unwind: a trader with dry capital buys the basis they sell.
\sfield{Instruments and venues} As for the basis long.
\sfield{Signal} Leveraged funds' short futures at extremes, collateral terms rising, the basis widening fast.
\sfield{Sizing and execution} Small, added as forced selling peaks; funded with term repo.
\sfield{Costs} Wide markets in stress.
\sfield{How it dies} Buying too early into a longer unwind.
\sfield{Horizon, capacity, infrastructure} Weeks; committed funding.
\sfield{Backtest honestly} Few episodes; the CFTC category mixes strategies.
\sfield{Sources} CFTC Traders in Financial Futures (this chapter's derived series); no performance figure verified.
\end{strategyfile}

\section{Tutorial: fifty times}\label{tut:s2:the-bond-futures-basis-trade:tut}

\begin{tutorial}
\textbf{Goal.} Run the basis book at several leverages through a funding stress, and read leveraged funds' positions in the CFTC data.
\textbf{End state:} the table and the two figures.
\begin{tutsteps}
\item \textbf{The book}.
\omcode{code/firm/basistrade/firm_basistrade.py}{66}{96}{Carry, marks, funding and forced sales, with re-levering at the rolls.}\label{lst:s2:the-bond-futures-basis-trade:book}
\item \textbf{By leverage}.
\omcode{code/strategies-2/11-the-bond-futures-basis-trade/python/s2_basis.py}{37}{50}{Return before, trough and loss after the stress, by leverage.}\label{lst:s2:the-bond-futures-basis-trade:table}
\item \textbf{Run} \texttt{s2\_fetch\_tff.py} once, then \texttt{table()}, \texttt{parts()} and \texttt{fig\_basis.py}.
\end{tutsteps}
\textbf{What to change next.} Finance with term repo that cannot be withdrawn during the stress and find the leverage it allows; hold a cash buffer
of 2\% and compare; let the stress come twice in a year.
\end{tutorial}

\section{Build: basis trade}\label{bld:s2:the-bond-futures-basis-trade:basistrade}

\begin{build}
\bfield{Purpose} A leveraged cash-futures basis book with haircuts, futures margin, repo cost and forced deleveraging.
\bfield{Interface} \texttt{BasisConfig(\dots)}, \texttt{simulate\_basis(cfg)}, \texttt{run\_book(sim, cfg, leverage)}.
\bfield{Rules} Capital covers haircut plus margin on the face held or the book sells; no re-levering in the stress; quarterly re-levering after.
\bfield{Acceptance tests} \texttt{code/firm/basistrade/tests/}: a calm book earns the spread times leverage; forced sales only when capital falls
short, and capital covers collateral afterwards; the stress scenario's shape.
\bfield{Stretch} Term repo; the delivery option; several contracts.
\end{build}

\begin{omsources}
\item D.~Barth and R.~J.~Kahn, ``Hedge funds and the Treasury cash-futures disconnect'', Office of Financial Research Working Paper 21-01, 2021.
\item Board of Governors of the Federal Reserve System, press releases of 15 and 23 March 2020.
\item CFTC, Traders in Financial Futures, 2010--2026.
\end{omsources}

\section{Exercises}

\begin{exercise}[$\star$]\label{exo:s2:the-bond-futures-basis-trade:1}
The haircut and the futures margin are 1\% each. What is the highest leverage? And at 3\% each?
\end{exercise}

\begin{exercise}[$\star$]\label{exo:s2:the-bond-futures-basis-trade:2}
A book earns 20 basis points a year on its face at 50 times its capital. What does it earn on capital, before costs?
\end{exercise}

\begin{exercise}[$\star$]\label{exo:s2:the-bond-futures-basis-trade:3}
The basis moves 60 cents per 100 against a book at 50 times. What does it lose on capital, before any sale?
\end{exercise}

\begin{exercise}[$\star\star$]\label{exo:s2:the-bond-futures-basis-trade:4}
Why does the ten-times book recover and the fifty-times book not?
\end{exercise}

\begin{exercise}[$\star\star$]\label{exo:s2:the-bond-futures-basis-trade:5}
By how much did hedge funds cut their short 2-, 5- and 10-year futures between 18 February and 17 March 2020, in percent?
\end{exercise}

\begin{exercise}[$\star\star$]\label{exo:s2:the-bond-futures-basis-trade:6}
Why is the CFTC's leveraged-funds series only a rough measure of the basis trade?
\end{exercise}

\begin{exercise}[$\star\star\star$]\label{exo:s2:the-bond-futures-basis-trade:7}
\emph{Coding.} Rerun the fifty-times book with \texttt{BasisConfig(haircut\_stress=0.01, margin\_stress=0.01)}, so that collateral terms do not
change. What happens, and what does it show?
\end{exercise}

\begin{exercise}[$\star\star\star$]\label{exo:s2:the-bond-futures-basis-trade:8}
\emph{Find the flaw.} ``The basis always converges at delivery, so the trade cannot lose if we hold to delivery.''
\end{exercise}

\section{Problem: Fifty Times}

\begin{problem}[{Weekend problem --- a leveraged basis book}]\label{pb:s2:the-bond-futures-basis-trade:1}
The chapter's synthetic book, the CFTC data and the public record.

\medskip\noindent\textbf{Part I --- Mechanics.}
\begin{enumerate}
\item Define the gross basis, net basis and implied repo rate.
\item When does a long basis position earn?
\item Define \omterm{def:s2:the-bond-futures-basis-trade:leverage}{basis-trade leverage} and \omterm{def:s2:the-bond-futures-basis-trade:rollover}{repo rollover risk}.
\item What caps the leverage?
\end{enumerate}

\medskip\noindent\textbf{Part II --- Who does it.}
\begin{enumerate}[resume]
\item Why do asset managers push futures above fair value?
\item What did Barth and Kahn estimate about size?
\item What does the CFTC series show?
\item Why do the authors call hedge funds warehouses?
\end{enumerate}

\medskip\noindent\textbf{Part III --- March 2020.}
\begin{enumerate}[resume]
\item What happened in the Treasury and repo markets?
\item What did hedge funds do?
\item What did the Federal Reserve announce?
\item Were basis trades the primary cause, according to Barth and Kahn?
\end{enumerate}

\medskip\noindent\textbf{Part IV --- The verdict.}
\begin{enumerate}[resume]
\item State the \emph{named result}: the trade's return on capital by leverage and the loss in the planted funding stress.
\item Decompose the fifty-times book's loss.
\item What would term funding change?
\item How would you measure crowding?
\item How would you size the book?
\item How would you backtest it honestly?
\item Which strategy file profits from others' forced sales?
\item In one sentence: what does a basis trader provide to the market?
\end{enumerate}
\end{problem}

\section{Interview questions}

\begin{interviewq}[$\star$ \iqroles{trader}]\label{iq:s2:the-bond-futures-basis-trade:1}
Explain the Treasury cash-futures basis trade.
\end{interviewq}

\begin{interviewq}[$\star\star$ \iqroles{risk}]\label{iq:s2:the-bond-futures-basis-trade:2}
A basis book is levered 50 times. What stress would you run, and what limit would you set?
\end{interviewq}

\begin{interviewq}[$\star\star$ \iqroles{researcher}]\label{iq:s2:the-bond-futures-basis-trade:3}
How would you compute the implied repo rate, and what does its gap to repo tell you?
\end{interviewq}

\begin{interviewq}[$\star\star$ \iqroles{trader}]\label{iq:s2:the-bond-futures-basis-trade:4}
Your repo lender raises the haircut from 1\% to 3\% overnight. What do you do?
\end{interviewq}

\begin{interviewq}[$\star\star$ \iqroles{developer}]\label{iq:s2:the-bond-futures-basis-trade:5}
Design the daily process for a basis book: data, collateral, margin calls, rolls.
\end{interviewq}

\begin{interviewq}[$\star\star\star$ \iqroles{researcher}]\label{iq:s2:the-bond-futures-basis-trade:6}
Show that, with capital $C$, face $N$ and collateral rate $k$ per unit of face, a loss of $x$ per unit of face forces sales when $C - Nx < Nk$, and find
the largest move a book at leverage $L$ can take without selling.
\end{interviewq}
